\chapter{总结与展望}
\label{chap:conclusion}

\section{研究成果总结}
\label{sec:achievement}

本课题以高级驾驶辅助系统在雨天环境下面临的车道线检测鲁棒性不足为核心问题，围绕“雨天干扰机理分析—专用数据集构建—轻量化抗干扰算法设计—嵌入式平台实时部署”的技术主线，开展系统性研究工作，取得了以下四个方面的主要研究成果。

\subsection{雨天车道线视觉特征退化机理的系统分析}

针对雨滴遮挡、水膜镜面反射及大气散射三类主要干扰因素，本文从物理成像模型与视觉特征两个层面展开了定量分析与实验验证。在物理模型层面，建立了雨滴遮挡区域的像素混合模型（式\ref{eq:raindrop_occlusion}），揭示了附着雨滴导致局部图像信息缺失的作用规律；引入菲涅耳反射定律描述了水膜覆盖路面的镜面反射效应（式\ref{eq:fresnel}），阐明了高亮伪边缘的生成机制；并基于Koschmieder大气散射模型（式\ref{eq:koschmieder}）刻画了降雨场景下全局对比度随场景深度指数衰减的退化趋势。在视觉特征层面，通过统计不同降雨强度下车道线区域的Michelson对比度、Canny边缘响应强度及局部遮挡率三项定量指标，揭示了车道线可辨识性随降雨强度增加而显著退化的统计规律——大雨条件下对比度较晴天下降约68\%，局部遮挡率可达30\%以上。进一步地，以CLRNet为基线模型在雨天模拟场景下进行失效模式归因分析，明确了水膜反射假阳性误检（占误检总数38\%）、雨滴遮挡断点漏检（占漏检总数45\%）及雨纹方向混淆误检（占误检总数17\%）三类典型失效模式。该部分工作为后续数据集构建策略的制定与抗干扰模块的设计提供了明确的理论导向与数据支撑。

\subsection{RainLane雨天车道线专用数据集的构建与系统增强}

针对现有公开数据集在雨天场景覆盖严重不足的结构性缺陷，本课题构建了面向雨天车道线检测的专用数据集RainLane。在数据采集环节，遵循多维度覆盖原则，于真实道路环境中系统采集了覆盖小雨、中雨、大雨/暴雨四级降雨强度、白天/夜间/傍晚三类光照条件、城市主干道/次干道/郊区公路/高速公路四类道路类型、以及路面湿润/积水/雨后潮湿三种路面状态的原始视频数据，经关键帧抽取与质量筛选，获得约28,000张有效原始图像。在数据标注环节，制定了适用于雨天模糊视距的专用标注规范，重点明确了基于车道线走向连续性的断点续标规则与不确定区域降权处理机制，并建立了“初级审核—一致性检验—模型辅助复审”三级质控流程，最终完成25,600张图像、98,400条车道线实例的高质量标注，数据集按8:1:1比例严格划分为训练集、验证集与测试集（详见表\ref{tab:rainlane_distribution}）。在数据增强环节，设计并实现了包含雨纹合成增强（式\ref{eq:rainstreak_synthesis}）、水膜反射模拟（式\ref{eq:reflection_simulation}）、CutOut遮挡增强及通用几何/光度增强的多层次增强流水线（算法\ref{alg:aug_pipeline}），将训练集有效样本扩充约4倍至81,920张，有效缓解了特定极端场景下的样本稀缺问题。跨数据集迁移实验表明，仅使用TuSimple晴天数据训练的CLRNet模型在RainLane测试集上F1分数仅为43.2\%，而在RainLane训练集上微调后F1分数提升至76.8\%，验证了本数据集对提升雨天检测性能的关键作用。

\subsection{面向雨天环境的轻量化抗干扰车道线检测算法设计}

针对基线模型CLRNet在雨天场景下特征中断、伪边缘误激活及行分类模糊等失效问题，本文提出了一种即插即用的频域门控注意力模块，并将其嵌入CLRNet主干网络与跨层精炼模块之间，构建了CLRNet-FGM模型。FGM模块的核心设计思路源于对雨天图像干扰频域特性的分析——雨纹噪声能量集中于高频方向性尖峰，而车道线边缘信息分布于中高频定向能量带。基于此，模块设计了三个子单元：\textbf{多尺度频域分解单元}采用分组深度可分离卷积（式\ref{eq:multiscale_freq}），通过不同尺寸的卷积核近似模拟不同感受野下的频率响应，实现对高频雨纹噪声与中频水膜反射干扰的分级捕获；\textbf{门控权重生成单元}通过全局平均池化聚合频域响应统计量，经两层全连接网络输出通道级门控向量（式\ref{eq:gate_generation}），赋予模型自适应调控不同频率分量响应的能力；\textbf{特征重标定单元}利用门控向量对原始特征进行逐通道加权，并辅以残差连接（式\ref{eq:residual_recalibration}），在抑制干扰的同时保留有用信息的恒等传递路径。经计算，FGM模块的引入仅使模型参数量增加约4.2\%、FLOPs增加约3.8\%，满足轻量化设计约束。此外，本文设计了包含车道线位置损失（式\ref{eq:loss_pos}）、存在性分类损失（式\ref{eq:loss_cls}）及雨天鲁棒正则化项（式\ref{eq:loss_reg}）的综合损失函数，其中正则化项通过约束雨天图像与清晰图像的门控向量一致性，促使FGM模块学习对雨天干扰具有判别性的域不变特征。训练采用两阶段策略——先在CULane数据集预训练主干网络，后在RainLane数据集联合微调FGM模块与预测头，有效避免了从头训练时的收敛困难问题。

\subsection{嵌入式平台部署优化与实时性能验证}

为验证算法在真实ADAS车载平台上的工程化可行性，本文系统完成了从PyTorch模型到NVIDIA Orin嵌入式平台TensorRT推理引擎的端到端部署与优化。首先，将训练好的CLRNet-FGM模型导出为ONNX中间格式，并通过常量折叠、算子融合等图优化手段将计算图节点数量压缩约20\%。其次，基于TensorRT构建了支持FP16与INT8两种精度模式的推理引擎，采用熵校准算法对INT8模式进行训练后量化校准，校准数据集覆盖RainLane训练集中不同降雨强度与光照条件的500帧代表性样本。再次，在Orin平台上实现了双缓冲异步推理流水线与GPU-DLA混合推理模式，将主干网络的部分卷积层卸载至深度学习加速器执行，在保证功能正确性的同时使GPU占用率降低约25\%、整体功耗下降约15\%。实时性能测试结果（表\ref{tab:deploy_performance}）表明：在FP16精度下，模型推理延迟为15.8 ms，对应吞吐量63.3 FPS，远超ADAS系统30 FPS的实时性门槛；在INT8精度下，延迟进一步压缩至9.2 ms，吞吐量突破100 FPS，而检测精度较FP32基准仅下降约1.8个百分点。综合精度-速度评估显示，CLRNet-FGM以8.2\%的推理延迟增量为代价，换取了9.2个百分点F1分数的绝对提升，实现了鲁棒性与实时性的良好平衡。

\subsection{全面的实验评估与性能分析}

在RainLane测试集上对CLRNet-FGM模型进行了多维度实验验证。整体精度对比（表\ref{tab:overall_performance}）显示，CLRNet-FGM的F1分数达81.7\%，较基线CLRNet提升9.2个百分点，精确率与召回率分别提升8.2和10.0个百分点，在所有对比方法中排名第一。分降雨强度评估（表\ref{tab:rain_intensity}）表明，在大雨/暴雨场景下CLRNet-FGM的F1分数为76.8\%，相对基线提升11.1个百分点，且性能衰减幅度显著小于基线模型；雨天退化鲁棒性指数RDRI达0.88，较基线（0.67）提升21个百分点，证明模型对域偏移具有强抵抗力。消融实验结果（表\ref{tab:ablation}）揭示，FGM模块整体贡献了9.2个百分点的性能增益，其中多尺度频域分解设计与雨天鲁棒正则化项分别贡献3.5和2.2个百分点，验证了各组件的有效性。跨数据集泛化性测试显示，模型在BDD100K独立雨天子集上仍取得79.3\%的F1分数，较基线提升9.5个百分点，证明其学到的抗干扰能力具备良好的可迁移性。极端干扰鲁棒性测试进一步表明，在重度雨纹合成与随机块遮挡工况下，CLRNet-FGM的F1分数分别保持70.5\%和75.3\%，显著优于基线模型的54.2\%和61.8\%，展现了应对真实行车环境中不确定性干扰的潜力。

综上所述，本课题以“数据-算法-部署”三位一体的研究范式，系统性地解决了雨天条件下车道线检测鲁棒性不足与实时性难以兼顾的关键难题，形成了一套完整、可用、高效的解决方案，为高级驾驶辅助系统在恶劣天气下的安全可靠运行提供了有力的理论依据与技术支撑。

\section{创新点总结}
\label{sec:innovation}

本课题在已有研究基础上，针对雨天环境下车道线检测的特定挑战，在数据集构建方法、算法模型设计以及工程部署策略三个层面进行了创新性探索，主要创新点归纳如下。

\subsection{构建了首个系统标注降雨强度与路面状态的雨天车道线专用数据集}

现有公开数据集（如TuSimple、CULane）中雨天样本占比极低且缺乏精细化场景分级标签，严重制约了雨天车道线检测算法的针对性训练与公平评估。本课题首次从“降雨强度—光照条件—道路类型—路面状态”四个维度系统性构建了RainLane数据集，包含25,600张高质量标注图像，并创新性地引入了降雨强度等级（小雨/中雨/大雨/暴雨）、路面状态（湿润/积水/潮湿）等细粒度属性标签，为域自适应学习与分场景性能评估提供了数据基础。同时，所提出的雨天专用标注规范（含断点续标规则与三级质控机制）为后续类似数据集的构建提供了可参考的范式。该数据集的构建填补了领域内雨天标准化评测基准的空白，是本文工作的重要基础性贡献。

\subsection{提出了基于频域门控注意力的轻量化抗干扰特征增强机制}

针对传统图像去雨预处理与检测任务串联导致的特征失配与计算冗余问题，本文首次从频域分析视角切入，提出了一种可嵌入主干网络内部的频域门控注意力模块。FGM模块的创新之处体现在：其一，利用分组深度可分离卷积在特征通道维度上隐式实现多尺度频域分解，避免了显式频域变换（如FFT）带来的高昂计算开销，实现了频域处理与空间特征提取的解耦；其二，通过可学习的门控权重生成网络，赋予模型自适应调控不同频率分量响应的能力，使其能够根据输入图像的内容动态抑制雨纹高频噪声而增强车道线中频结构响应，相较于固定滤波器具有更强的场景适应性；其三，以即插即用的轻量化设计嵌入CLRNet主干网络，参数量与计算量增量均控制在5\%以内，在不显著牺牲实时性的前提下大幅提升了雨天检测鲁棒性。这一设计思路为解决低信噪比、强干扰场景下的细长目标检测问题提供了一种新的技术途径。

\subsection{设计了面向雨天域适应的正则化训练策略}

为强化FGM模块对雨天干扰的判别能力，本文在损失函数中引入了一项雨天鲁棒正则化项\(\mathcal{L}_{\mathrm{reg}}\)（式\ref{eq:loss_reg}）。该正则项以对比学习为思想内核，约束同一场景在雨天与清晰条件下的FGM门控向量趋于一致，迫使模块学习与天气条件无关的、鲁棒的车道线特征表达。消融实验证实，该正则项独立贡献了2.2个百分点的F1分数提升，表明其在引导模型学习域不变特征方面的有效性。这一策略为深度学习模型在恶劣天气场景下的域适应训练提供了简单而有效的正则化手段。

\subsection{实现了算法精度与边缘端实时性的协同优化与验证}

多数针对恶劣天气的鲁棒检测算法因计算复杂度高而难以满足车载嵌入式的实时性约束，本文从算法设计与工程部署两个层面协同优化了精度与速度的平衡。算法层面，FGM模块在设计之初即遵循轻量化原则，采用深度可分离卷积与通道注意力瓶颈结构，从根本上控制了计算增量。部署层面，系统探索了TensorRT FP16/INT8量化、GPU-DLA异构计算、双缓冲异步流水线等优化技术的组合应用，最终在NVIDIA Orin平台上以63.3 FPS的实时帧率稳定运行，且INT8模式下进一步突破100 FPS。这一“算法-部署”协同优化的实践路径，为同类视觉感知模型在资源受限边缘端的落地提供了可复现的工程范式与技术参考。

\section{未来工作展望}
\label{sec:future_work}

本课题在雨天车道线鲁棒检测方法研究方面取得了一定进展，但受限于研究时间与实验条件，仍存在若干问题有待进一步深入探索。未来工作可从以下四个方向继续推进。

\subsection{多传感器融合增强全天候感知鲁棒性}

本课题方法仅依赖单目视觉传感器，在极端恶劣天气（如特大暴雨导致能见度骤降、镜头被完全水膜覆盖）下，视觉信息严重退化，单纯依靠图像处理与深度学习模型的性能提升空间有限。未来研究可引入毫米波雷达与激光雷达等多源传感器数据，利用雷达点云在恶劣天气下仍能提供稳定几何结构信息的特点，与视觉特征进行多模态融合。具体而言，可探索将雷达点云投影至图像平面生成稀疏深度图作为注意力先验，引导视觉分支更加关注路面区域的潜在车道线位置；或设计基于Transformer的跨模态融合模块，在特征层面实现视觉纹理与雷达几何信息的自适应对齐与互补。多传感器融合有望从根本上突破单一视觉模态在极端天气下的性能瓶颈，是实现L3及以上级别自动驾驶全天候安全运行的必要技术路径。

\subsection{时序信息建模与动态雨天场景适应}

本课题方法以单帧图像为输入，未充分利用视频序列中蕴含的时序上下文信息。实际行车过程中，车道线的几何形态具有强时序连续性，雨滴的运动轨迹、水膜的动态变化也呈现出可预测的时空规律。未来工作可将当前帧级模型扩展至视频序列输入，引入ConvLSTM、3D卷积或时序Transformer等结构，显式建模前后帧之间的特征演化关系。时序建模可带来双重收益：一方面，利用历史帧的车道线检测结果对当前帧进行平滑约束，可有效抑制因瞬时雨滴遮挡或反光造成的检测跳变，提升输出结果的稳定性与平滑性；另一方面，可通过学习雨纹的动态模式，更精准地区分运动雨滴噪声与静态车道线结构，进一步提升抗干扰能力。此外，基于视频序列的时序一致性也可作为自监督信号，用于无标注雨天数据上的域适应训练，降低对高成本人工标注的依赖。

\subsection{面向更广泛恶劣天气场景的泛化能力研究}

本课题聚焦于雨天这一特定恶劣天气场景，但实际ADAS系统还需应对雾天、雪天、沙尘暴、强逆光等复杂环境。这些场景下车道线视觉特征的退化机理各有不同——雾天主要表现为全局对比度下降与细节模糊，雪天则面临路面积雪覆盖与雪花动态遮挡的双重挑战。FGM模块的频域门控思想在理论上具有一定通用性，未来可探索将其推广至多天气场景的统一框架中。一种可行的思路是构建天气条件编码向量作为FGM模块的额外输入条件，使模型能够根据当前天气类型自适应地调整频域滤波策略；或引入元学习机制，训练一个能够快速适应新天气域的初始化模型，仅需少量目标域样本即可实现有效的域迁移。构建覆盖雨、雾、雪、夜间等多场景的统一车道线检测模型，是推动ADAS环境感知模块向通用鲁棒性演进的重要研究方向。

\subsection{模型进一步轻量化与车规级芯片适配}

尽管CLRNet-FGM模型已在NVIDIA Orin平台上达到了63.3 FPS的实时推理速度，但距离部分低功耗车规级芯片（如TI TDA4VM、地平线征程系列）的算力约束仍有进一步压缩空间。未来可从以下子方向继续轻量化探索：其一，采用神经网络架构搜索技术在更大的搜索空间中自动发现针对雨天场景最优的轻量化网络结构，以替代人工设计的ResNet-18主干与FGM模块组合；其二，探索知识蒸馏策略，以当前高性能CLRNet-FGM模型为教师网络，指导参数量更小的学生网络（如MobileNetV4、EfficientNet-B0）学习雨天场景下的鲁棒特征表达；其三，针对特定车规级芯片的硬件加速器特性（如地平线BPU的指令集约束），进行算子级定制优化，最大限度发挥硬件算力潜能。此外，模型的功耗评估与热稳定性测试也是走向量产装车前的必要环节，有待后续补充开展。

\subsection{端到端自动驾驶系统闭环验证}

本课题仅在离线数据集上评估了车道线检测模块的独立性能，尚未将模型集成至完整的ADAS决策控制闭环中进行实车验证。在真实行车环境中，车道线检测结果将作为车道保持辅助、自动变道等下游功能的输入，检测误差在闭环系统中可能被累积放大。未来工作需将CLRNet-FGM模型部署至实车或高保真硬件在环仿真平台，在动态雨天场景下测试其对车辆横向控制精度、驾乘舒适性及任务成功率的影响。通过系统级闭环评估，可进一步发现检测算法在真实时序任务中的潜在缺陷，并为算法改进提供更具工程意义的反馈信号。

综上所述，本课题在雨天车道线鲁棒检测的数据、算法与部署三个层面取得了阶段性成果，但距离实现全天候、全场景可靠运行的高级驾驶辅助系统仍有持续研究的空间。未来将在多传感器融合、时序建模、多天气泛化、极致轻量化及系统闭环验证等方向继续深入探索，为提升自动驾驶技术在复杂环境下的安全性与适应性贡献力量。